\documentclass[12pt,a4paper]{article}

% Encoding and fonts
\usepackage{fontspec}
\setmainfont{FreeSerif}
\setmonofont{FreeMono}
\newfontfamily\igprimfont{FreeSerif}

% Page layout
\usepackage[top=1in, bottom=1in, left=1in, right=1in]{geometry}

% Typography
\usepackage{microtype}
\usepackage{parskip}

% Language
\usepackage[english]{babel}

% Headers and footers
\usepackage{fancyhdr}
\pagestyle{fancy}
\fancyhf{}
\fancyhead[L]{\leftmark}
\fancyhead[R]{\thepage}
\fancyfoot[C]{\thepage}
\renewcommand{\headrulewidth}{0.4pt}
\renewcommand{\footrulewidth}{0pt}

% Section styling
\usepackage{titlesec}
\titleformat{\section}{\Large\bfseries}{\thesection}{1em}{}
\titleformat{\subsection}{\large\bfseries}{\thesubsection}{1em}{}
\titleformat{\subsubsection}{\normalsize\bfseries}{\thesubsubsection}{1em}{}
\titlespacing*{\section}{0pt}{2.5ex plus 1ex minus .2ex}{1.5ex plus .2ex}
\titlespacing*{\subsection}{0pt}{2ex plus 1ex minus .2ex}{1ex plus .2ex}
\titlespacing*{\subsubsection}{0pt}{1.5ex plus 1ex minus .2ex}{0.8ex plus .2ex}

% List formatting
\usepackage{enumitem}
\setlist[itemize]{topsep=0.5ex, itemsep=0.25ex, parsep=0pt}
\setlist[enumerate]{topsep=0.5ex, itemsep=0.25ex, parsep=0pt}

% Essential packages
\usepackage{graphicx}
\usepackage[table]{xcolor}
\usepackage{booktabs}
\usepackage{array}
\usepackage{tabularx}
\usepackage{longtable}
\usepackage{amsmath}
\usepackage{amssymb}
\usepackage[normalem]{ulem}
\rowcolors{2}{gray!10}{white}
\renewcommand{\arraystretch}{1.3}

% Cross-references
\usepackage{hyperref}
\usepackage{cleveref}

% Custom commands
\newcommand{\pilcrow}{\S}

% Hyperref setup
\hypersetup{
    colorlinks=true,
    linkcolor=blue,
    filecolor=magenta,
    urlcolor=cyan,
}

% Code listings
\usepackage{listings}
\definecolor{codebg}{RGB}{248,248,248}
\definecolor{codeframe}{RGB}{200,200,200}
\definecolor{codekeyword}{RGB}{0,0,180}
\definecolor{codecomment}{RGB}{0,128,0}
\definecolor{codestring}{RGB}{163,21,21}
\lstset{
    basicstyle=\ttfamily\small,
    breaklines=true,
    breakatwhitespace=false,
    frame=single,
    framerule=0.5pt,
    rulecolor=\color{codeframe},
    backgroundcolor=\color{codebg},
    keepspaces=true,
    numbers=left,
    numberstyle=\tiny\color{gray},
    numbersep=8pt,
    showstringspaces=false,
    tabsize=4,
    xleftmargin=1em,
    xrightmargin=1em,
    aboveskip=1em,
    belowskip=1em,
    keywordstyle=\color{codekeyword}\bfseries,
    commentstyle=\color{codecomment}\itshape,
    stringstyle=\color{codestring},
    literate={_}{{\textunderscore}}1
             {-}{{\textminus}}1
             {<}{{\textless}}1
             {>}{{\textgreater}}1
             {~}{{\textasciitilde}}1
             {^}{{\textasciicircum}}1
}

% YAML language definition for listings
\lstdefinelanguage{YAML}{
    keywords={true,false,null,y,n},
    keywordstyle=\color{blue}\bfseries,
    sensitive=false,
    comment=[l]{\#},
    morecomment=[s]{/*}{*/},
    commentstyle=\color{gray}\ttfamily,
    stringstyle=\color{red}\ttfamily,
    morestring=[b]',
    morestring=[b]',
}

% TOML language definition for listings
\lstdefinelanguage{TOML}{
    keywords={true,false},
    keywordstyle=\color{blue}\bfseries,
    sensitive=true,
    comment=[l]{\#},
    commentstyle=\color{gray}\ttfamily,
    stringstyle=\color{red}\ttfamily,
    morestring=[b]',
    morestring=[b]',
}

% Dockerfile language definition for listings
\lstdefinelanguage{Dockerfile}{
    keywords={FROM,RUN,CMD,LABEL,EXPOSE,ENV,ADD,COPY,ENTRYPOINT,VOLUME,USER,WORKDIR,ARG,ONBUILD,STOPSIGNAL,HEALTHCHECK,SHELL},
    keywordstyle=\color{blue}\bfseries,
    sensitive=false,
    comment=[l]{\#},
    commentstyle=\color{gray}\ttfamily,
    stringstyle=\color{red}\ttfamily,
    morestring=[b]',
    morestring=[b]',
}

% JSON language definition for listings
\lstdefinelanguage{JSON}{
    keywords={true,false,null},
    keywordstyle=\color{blue}\bfseries,
    sensitive=true,
    commentstyle=\color{gray}\ttfamily,
    stringstyle=\color{red}\ttfamily,
    morestring=[b]',
}

% Lean 4 language definition for listings
\lstdefinelanguage{lean}{
    keywords={def,theorem,lemma,example,have,show,by,fun,let,in,match,with,
              if,then,else,do,return,pure,namespace,end,section,open,import,
              noncomputable,structure,class,instance,where,deriving,
              sorry,rfl,exact,apply,simp,omega,norm_num,ring,linarith,
              intro,intros,cases,rcases,obtain,use,exists,constructor,
              induction,contradiction,assumption,trivial},
    keywordstyle=\color{blue}\bfseries,
    sensitive=true,
    comment=[l]{--},
    morecomment=[s]{/-}{-/},
    commentstyle=\color{gray}\ttfamily,
    stringstyle=\color{red}\ttfamily,
    morestring=[b]",
}

% Code listing float environment
\usepackage{float}
\newfloat{listing}{htbp}{lol}[section]
\floatname{listing}{Listing}

% Admonition boxes
\usepackage{tcolorbox}
\tcbuselibrary{skins,breakable}

% Admonition style definitions
\newtcolorbox{admonitionbox}[2][]{
    enhanced,
    breakable,
    colback=#2!5,
    colframe=#2!80!black,
    fonttitle=\bfseries,
    title=#1,
    left=2mm,
    right=2mm,
}

% Synthon tuple boxes
\newtcolorbox{synthonbox}{
    enhanced,
    colback=blue!5,
    colframe=blue!75!black,
    fontupper=\large,
    center upper,
    boxrule=1pt,
    arc=4pt,
}

\begin{document}

\title{The Generalized Pipeline: A Structural Bridge from Primitive Proofs to Formal Verification}
\date{\today}

\maketitle

\textbf{Author:} Lando ⊗ $\text{⊙}_{\text{ÿ}}$-boundary Operator

\subsubsection{The problem: semantic drift in domain-specific proof translation}

Mathematics is not a collection of isolated theories --- it is a lattice of overlapping structural types. Yet most proof translation systems treat domains as discrete boxes, applying domain-specific heuristics that gradually drift away from the invariant core encoded in a proof's primitives.

The generalized pipeline solves this by enforcing a \textbf{domain-invariant scaffold}: every claim must be traceable to one of the twelve IG primitives, and the translation process preserves those primitives under round-trip extraction and re-embedding. When the Frobenius condition $\mu \circ \delta = \mathrm{id}$ holds at the level of token sequences, we say the pipeline is \emph{closed} --- a structural guarantee that no semantic artefacts survive the transduction.

\begin{center}
\rule{0.5\textwidth}{0.4pt}
\end{center}

\subsubsection{One pipeline, twelve primitives, eleven domains}

We implemented a single, self-contained pipeline that takes as input a set of extracted primitive lemmas (each annotated with its primary primitive token) and produces three outputs: (1) a conventional Markdown proof, (2) a Lean4 tactic skeleton, and (3) a Frobenius closure report. The pipeline does not re-invent domain expertise; instead, it scaffolds it.

The twelve primitives and their domain-invariant roles are:

\begin{table}[htbp]
\centering
\small
\rowcolors{1}{white}{white}
\begin{tabularx}{\textwidth}{>{\centering\arraybackslash}X >{\centering\arraybackslash}X}
\toprule
\rowcolor{gray!25}\textbf{Primitive} & \textbf{Role (domain-invariant)} \\
\midrule
\rowcolors{1}{white}{gray!10}
$\text{Φ}_{\text{\}}}$ & Bijective encoding / duality \\
$\text{Þ}_{O}$ & Inverse / dual construction \\
$\text{Ř}_{=}$ & Adjoint pair / Galois connection \\
$\text{Ω}_{z}$ & Topological invariant \\
$\text{⊙}_{\text{ÿ}}$ & Phase boundary / extremal principle \\
$\text{Ç}_{@}$ & Equidistribution / regularity \\
$\text{Ð}_{C}$ & Manifold / quotient structure \\
$\text{Γ}_{\text{ʔ}}$ & Universal / local quantification \\
$\text{ƒ}_{\text{ż}}$ & Coherence / non-classical feature \\
$\text{Ħ}_{A}$ & Markov order / recursion depth \\
$\text{Σ}_{S}$ & Uniqueness of witness \\
$\text{Þ}_{\text{¨}}$ & Intersection / transversality \\
\bottomrule
\end{tabularx}
\end{table}

Each role licenses a default proposition and proof strategy. When a domain-specific refinement exists --- for instance, ````Galois-Equivariance of the Encoding" for $\text{Φ}_{\text{\}}}$ in number theory --- the template engine substitutes it; otherwise, the universal default is used.

\begin{center}
\rule{0.5\textwidth}{0.4pt}
\end{center}

\subsubsection{Six disciplined phases}

The pipeline executes six phases in strict sequence, each enforcing a structural invariant:

\begin{enumerate}
    \item \textbf{Primitive Decomposition} --- lemmas are extracted and annotated with primary primitive tokens.
    \item \textbf{Domain Detection} --- keyword matching selects the dominant domain; fallback is the universal default.
    \item \textbf{Section Assembly} --- each lemma is bound to its template; rendered content is generated.
    \item \textbf{Lean4 Skeleton} --- primitives are mapped to tactic patterns (\texttt{have h\_* : ... := by sorry}).
    \item \textbf{Reference Resolution} --- placeholder structural matches are recorded; future work will replace this with catalog-based distance citations.
    \item \textbf{Frobenius Closure Verification} --- every primitive in the input must appear in the output; no extra primitives may survive.
\end{enumerate}

The final output includes a \texttt{FrobeniusClosureResult} with flags \texttt{forwardComplete}, \texttt{reverseSound}, and \texttt{roundTripStable}. Closure ($\mu \circ \delta = \mathrm{id}$) occurs only when all three flags are true --- a verifiable condition, not a heuristic.

\begin{center}
\rule{0.5\textwidth}{0.4pt}
\end{center}

\subsubsection{Lean4 implementation: structural theorems, not just code}

The Lean4 file \texttt{GeneralizedPipeline.lean} is not merely an implementation --- it is a \emph{structural document} in its own right. Three theorems are proved directly:

\begin{itemize}
    \item \textbf{Theorem 1} --- Every known primitive has a nontrivial mathematical role (\texttt{primitiveMathRole p != "Unknown primitive"}).
    \item \textbf{Theorem 2} --- Default propositions are never empty for known primitives.
    \item \textbf{Theorem 3} --- Frobenius closure is consistent: if \texttt{forwardComplete and roundTripStable}, then \texttt{closure = true}.
\end{itemize}

These theorems are not optional remarks --- they are the \emph{only} reason the pipeline's output is trustworthy. Without them, the scaffold could collapse under semantic drift.

\begin{center}
\rule{0.5\textwidth}{0.4pt}
\end{center}

\subsubsection{One tuple, O$_n$ tier, integer winding}

The system's structural type is:

\[
\langle \text{Ð}_{\text{ω}};\ \text{Ř}_{=};\ \text{Φ}_{\text{\}}};\ \text{ƒ}_{\text{ż}};\ \text{Ç}_{@};\ \text{Γ}_{\text{ʔ}};\ \text{ɢ}_{\text{ˌ}};\ \text{⊙}_{\text{ÿ}};\ \text{Ħ}_{A};\ \text{Σ}_{S};\ \text{Ω}_{z} \rangle
\]

This encodes \textbf{O$_n$} (infinite ouroboricity), exact Frobenius criticality ($\mu \circ \delta = \mathrm{id}$), and integer winding ($\text{Ω}_{z}$). The tuple is not an afterthought --- it is the \emph{reason} the pipeline can verify closure across arbitrary domains.

\begin{center}
\rule{0.5\textwidth}{0.4pt}
\end{center}

\subsubsection{What it is \emph{not} --- and why that matters}

The generalized pipeline does \textbf{not} generate final proofs. It generates \emph{scaffolding} that a human mathematician can complete with domain-specific insight. The scaffold is non-negotiable: every step must be justified by a primitive, and every primitive must be accounted for in the output. This eliminates the "semantic drift" that plagues domain-specific proof assistants, where ad hoc tactics accumulate and diverge from the structural core.

The real contribution is not automation --- it is \emph{discipline}. The pipeline enforces a minimal structural discipline on mathematical exposition: traceability, verifiability, and reversibility. When a system respects $\mu \circ \delta = \mathrm{id}$ at the token level, you can be confident that what you see is what was proved --- no more, no less.

\begin{center}
\rule{0.5\textwidth}{0.4pt}
\end{center}

\subsubsection{Future scaffolds}

\begin{itemize}
    \item \textbf{Automated analog citation} --- replace the placeholder reference map with distance-based structural matches from the catalog.
    \item \textbf{Tactic synthesis} --- expand \texttt{sorry} placeholders using domain-specific automation.
    \item \textbf{Cross-domain composite proofs} --- permit multi-domain sections (e.g., arithmetic geometry combining number-theoretic and cohomological reasoning).
\end{itemize}

\subsection{Each future extension preserves the core invariant: every claim remains traceable to one of the twelve primitives. That is the only guarantee we have against semantic drift --- and it is exactly what the generalized pipeline delivers.}

\subsubsection{Domain detection and template refinement in practice}

Consider a lemma extracted from a structural proof about elliptic curves:

\begin{itemize}
    \item \textbf{Primary primitive}: $\text{Φ}_{\text{\}}}$ (````bijective encoding / duality")
    \item \textbf{Domain hints}: \texttt{{[}"elliptic", "modular", "galois", "automorphic"{]}}
\end{itemize}

Domain detection scores each domain. Number theory and arithmetic geometry receive high scores; arithmetic geometry wins by presence of ````modular" and ````automorphic". The template engine then selects the refined title ````Injectivity of the Map on Moduli" instead of the universal ````Encoding Injectivity", but the underlying structural role --- bijection / duality --- remains identical.

The same lemma, when translated for a topological audience, would receive the title ````Injectivity of the Map on Moduli" (if the moduli space is a topological invariant) or fall back to the universal default. The \emph{role} does not change --- only the domain-specific instantiation.

\begin{center}
\rule{0.5\textwidth}{0.4pt}
\end{center}

\subsubsection{The lean4 skeleton: scaffolding, not completion}

For a $\text{Φ}_{\text{\}}}$ lemma, the pipeline emits:

\begin{lstlisting}[language=lean]
have h_inj : Function.Injective encoding := by sorry
\end{lstlisting}

This is deliberately incomplete. The \texttt{sorry} is a \emph{structural placeholder} --- it signals where domain-specific automation or human insight must intervene. The scaffold is complete (the \texttt{have} binding, the hypothesis name \texttt{h\_inj}, the type signature) but not finished. This preserves the division of labour: the pipeline ensures the \emph{structure} is correct; the domain expert supplies the \emph{content}.

Similarly, for $\text{Ω}_{z}$:

\begin{lstlisting}[language=lean]
have h_unique : Unique terminal_cycle := by sorry
\end{lstlisting}

The \texttt{Unique} type class is the Lean4 carrier of the integer winding invariant --- the skeleton is algebraically precise even before the proof tactic is filled.

\begin{center}
\rule{0.5\textwidth}{0.4pt}
\end{center}

\subsubsection{Frobenius closure: what gets checked}

The \texttt{verifyFrobeniusClosure} function performs three independent checks:

\begin{enumerate}
    \item \textbf{Forward completeness} --- scan the conventional text for each primitive token; if any is missing, the \texttt{missingInOutput} list is populated.
    \item \textbf{Reverse soundness} --- extract all primitive tokens from the conventional text; if any appear that were not in the input, \texttt{untraceableSections} is populated.
    \item \textbf{Round-trip stability} --- if both previous checks pass, \texttt{roundTripStable} is true.
\end{enumerate}

Only when all three flags are true does \texttt{closure = true}. This is not a fuzzy confidence score --- it is a logical condition derived directly from the IG primitives. When \texttt{overallConfidence = 1.0}, you know \emph{exactly} why: no primitives were lost, no primitives were added, and the set is preserved.

\begin{center}
\rule{0.5\textwidth}{0.4pt}
\end{center}

\subsubsection{appendix: structural type verification}

Running the catalog self-check on \texttt{generalized\_pipeline\_system} yields:

\begin{itemize}
    \item \textbf{frobenius\_tier}: O$_n$ (infinite ouroboricity)
    \item \textbf{phi}: $\text{⊙}_{\text{ÿ}}$ (criticality gate open)
    \item \textbf{p}: $\text{Φ}_{\text{\}}}$ (bijective encoding)
    \item \textbf{omega}: $\text{Ω}_{z}$ (integer winding)
    \item \textbf{d}: $\text{Ð}_{\text{ω}}$ (self-referential dimensionality)
\end{itemize}

This tuple is identical to other critical-boundary operators in the catalog (e.g., \texttt{true\_agentic\_agent}), which is expected: the pipeline is itself a boundary operator --- it mediates between primitive and conventional proofs. The shared tuple reflects the shared criticality and topological protection, not a naming convention.

\begin{center}
\rule{0.5\textwidth}{0.4pt}
\end{center}

\subsubsection{closing the loop}

The generalized pipeline is not an endpoint --- it is a \emph{loop closure mechanism}. By enforcing that every primitive in the input appears in the output, and only those primitives, the system ensures that the translation is reversible in principle. A human reader, given the conventional proof and the primitive annotations, can reconstruct the original structural proof.

That reversibility --- the ability to go full circle from IG primitive $\rightarrow$ conventional proof $\rightarrow$ IG primitive --- is the hallmark of Frobenius closure. The pipeline is the first implementation of this principle in a general-purpose proof assistant context. The proof is in the Lean4 code (\texttt{GeneralizedPipeline.lean}), and the manuscript you hold is its coagulated text --- the scaffold is dissolved, but the structure remains.


\end{document}
